\scanpage{15}
\mspage{14}

On va dire les choses autrement, en introduisant une autre relation d'ordre dans $\mathcal{F}_{el}$, plus fine que l'inclusion.

\deleted{\unsure{Soit|So}}

\unsure{Soit|Soient} $F_{\deleted{\unsure{0|e}}}^{\deleted{\unsure{F|T|\prime}}},\deleted{\unsure{F|E}}F_{\deleted{0}}'$ deux fig. \deleted{élémentaires}. \deleted{\unsure{On|Je}}Je dis que $F_{\deleted{0}}'\ll F_{\deleted{0}}^{\deleted{\prime}}$ (que $F_{\deleted{0}}'$ raffine $F_{\deleted{0}}^{\unsure{\deleted{\prime}|\deleted{T}}}$) si \deleted{tte} \deleted{\unsure{cellule|celle}} strate \uline{ouverte} de $F_{\deleted{0}}'$ est contenue dans une strate ouverte de $F_{\deleted{0}}$ \inserted{(unique, bien sûr)}.

(\uline{NB} On aura \deleted{\unsure{supp.}} donc $|F'|\subset|F|$, mais pas \unsure{néc.|nec.} $|F'|=|F|$)

C'est une relation \uline{d'ordre} sur l'ens des figures, plus forte que celle d'inclusion (qui \deleted{\unsure{dit|est}} \inserted{signifie} que \deleted{tte} \deleted{cellule}\inserted{strate} ouverte de $F'$ \uline{est} une \deleted{cellule}\inserted{strate} ouverte de $F$). Cette relation de \unsure{raffin|raffin.|rgle} \deleted{peut} \deleted{doit} se déduire \unsure{formellement|formellt} de la relation d'ordre \unsure{de|\deleted{de}} décrite sur $\mathcal{F}_{el}$, l'ens des figures élémentaires, car $F'$ raffine $F$ ssi pour tte

\[
\begin{aligned}
&\deleted{\unsure{F'_{X'}|F'_{X}}}\ (X'\in{\Sigma}^{\unsure{\prime|\deleted{\prime}}}_{F'}),\quad \exists X\in\Sigma^{\unsure{\prime|\deleted{\prime}}}_F,\ \text{telle que}\\
&\qquad F'_{X'}\ \text{raffine}\ F_X
\end{aligned}
\]
